\documentclass[10pt,twocolumn,letterpaper]{article}

% ---- layout in the style of the usual ML conference templates --------------
\usepackage[margin=0.75in,columnsep=0.25in]{geometry}
\usepackage{times}
\usepackage[T1]{fontenc}
\usepackage[utf8]{inputenc}
\usepackage{microtype}
\usepackage{amsmath,amssymb}
\usepackage{booktabs}
\usepackage{multirow}
\usepackage{array}
\usepackage[table]{xcolor}
\usepackage[numbers,sort&compress]{natbib}
\usepackage[colorlinks=true,linkcolor=blue!50!black,citecolor=blue!50!black,urlcolor=blue!50!black]{hyperref}
\usepackage[font=small,labelfont=bf,skip=4pt]{caption}
\usepackage{enumitem}
\setlist{nosep,leftmargin=*}
\usepackage{pgfplots}
\pgfplotsset{compat=1.18}
\usetikzlibrary{positioning,calc}
\usepackage{titlesec}
\titlespacing*{\section}{0pt}{8pt plus 2pt minus 2pt}{4pt plus 1pt}
\titlespacing*{\subsection}{0pt}{6pt plus 2pt minus 2pt}{3pt plus 1pt}
\titleformat{\section}{\normalfont\large\bfseries}{\thesection}{0.6em}{}
\titleformat{\subsection}{\normalfont\normalsize\bfseries}{\thesubsection}{0.6em}{}
\setlength{\parskip}{0pt}
\setlength{\textfloatsep}{8pt plus 2pt minus 2pt}

\setlength{\bibsep}{1pt plus 0.3ex}

\newcommand{\sys}{\textsc{SkillStore}}
\newcommand{\code}[1]{\texttt{\small #1}}

\newcommand{\NGPTmini}{39}
\newcommand{\FailGPTmini}{0}
\newcommand{\ValidZeroGPTmini}{56}
\newcommand{\ValidFinalGPTmini}{87}
\newcommand{\ValidOneGPTmini}{85}
\newcommand{\SecsGPTmini}{7}
\newcommand{\BlockedGPTmini}{5}
\newcommand{\NGPT}{38}
\newcommand{\FailGPT}{1}
\newcommand{\ValidZeroGPT}{89}
\newcommand{\ValidFinalGPT}{100}
\newcommand{\ValidOneGPT}{97}
\newcommand{\SecsGPT}{6}
\newcommand{\BlockedGPT}{0}
\newcommand{\NSonnet}{39}
\newcommand{\FailSonnet}{0}
\newcommand{\ValidZeroSonnet}{87}
\newcommand{\ValidFinalSonnet}{100}
\newcommand{\ValidOneSonnet}{100}
\newcommand{\SecsSonnet}{9}
\newcommand{\BlockedSonnet}{0}
\newcommand{\NOpus}{39}
\newcommand{\FailOpus}{0}
\newcommand{\ValidZeroOpus}{95}
\newcommand{\ValidFinalOpus}{100}
\newcommand{\ValidOneOpus}{100}
\newcommand{\SecsOpus}{24}
\newcommand{\BlockedOpus}{0}
\newcommand{\NAll}{155}
\newcommand{\BlockedFirstAll}{18}
\newcommand{\BlockedFinalAll}{3}
\newcommand{\BlockedFirstN}{28}
\newcommand{\BlockedFinalN}{5}
\newcommand{\UnmetFirst}{25}
\newcommand{\FlowFirst}{3}
\newcommand{\ObjectFirst}{1}
\newcommand{\IssuesFirst}{29}
\newcommand{\PointlessFirst}{12}
\newcommand{\PointlessFinal}{6}
\newcommand{\StepsTotal}{2010}
\newcommand{\ShareTrained}{87}
\newcommand{\ShareUntrained}{2}
\newcommand{\ShareGap}{11}
\newcommand{\GapTotal}{218}
\newcommand{\GapWithData}{99}
\newcommand{\DistinctGaps}{91}
\newcommand{\ApiFailures}{1}
\newcommand{\GapEgoShare}{75}
\newcommand{\NSkills}{59}
\newcommand{\NTrained}{11}
\newcommand{\NBeta}{5}
\newcommand{\NCatalogOnly}{43}
\newcommand{\NDatasets}{138}
\newcommand{\NOpenDatasets}{137}
\newcommand{\NLinks}{165}
\newcommand{\MedianEpisodes}{443}
\newcommand{\MinEpisodes}{50}
\newcommand{\TotalEpisodes}{317{,}112}
\newcommand{\NPredicates}{48}
\newcommand{\NKinds}{8}
\newcommand{\NActsOn}{8}
\newcommand{\Runs}{3}
\newcommand{\StabGPTmini}{0.86}
\newcommand{\SameGPTmini}{23}
\newcommand{\StabGPT}{0.93}
\newcommand{\SameGPT}{38}
\newcommand{\StabSonnet}{0.85}
\newcommand{\SameSonnet}{28}
\newcommand{\StabOpus}{0.90}
\newcommand{\SameOpus}{31}
\newcommand{\BalGPTmini}{1.45}
\newcommand{\SpeedGPTmini}{1.00}
\newcommand{\BothGPTmini}{18}
\newcommand{\BalGPT}{1.28}
\newcommand{\SpeedGPT}{1.00}
\newcommand{\BothGPT}{26}
\newcommand{\BalSonnet}{1.38}
\newcommand{\SpeedSonnet}{1.00}
\newcommand{\BothSonnet}{18}
\newcommand{\BalOpus}{1.38}
\newcommand{\SpeedOpus}{1.00}
\newcommand{\BothOpus}{36}
\newcommand{\SpeedAll}{1.00}
\newcommand{\SpeedMax}{2.0}
\newcommand{\BothAll}{25}
\newcommand{\SavedAll}{0}
\newcommand{\NScheduled}{150}
\newcommand{\ParShare}{25}
\newcommand{\ParShareBal}{62}
\newcommand{\SpeedMeanDefault}{1.13}
\newcommand{\SpeedMeanBal}{1.37}
\newcommand{\SpeedMaxBal}{2.0}
\newcommand{\TimeSavedBal}{29}
\newcommand{\DistinctVerbs}{27}
\newcommand{\MarketSkillsRaw}{99}
\newcommand{\MarketSkillsCanon}{35}
\newcommand{\MarketRequests}{99}
\newcommand{\MarketPacingRaw}{33}
\newcommand{\MarketPacingCanon}{45}
\newcommand{\TopBidRaw}{10.41}
\newcommand{\TopBidCanon}{20.21}
\newcommand{\NTeleop}{129}
\newcommand{\NEgo}{6}
\newcommand{\NSim}{3}
\def\TabModels{GPT-4.1 mini & 56 & 85 & \textbf{87} \scriptsize{[73,94]} & 95 & 11.5 & 0.8 & 7 \\
GPT-4.1 & 89 & 97 & \textbf{100} \scriptsize{[91,100]} & 100 & 10.8 & 1.7 & 6 \\
Claude Sonnet 5.5 & 87 & 100 & \textbf{100} \scriptsize{[91,100]} & 100 & 10.5 & 1.3 & 9 \\
Claude Opus 5.5 & 95 & 100 & \textbf{100} \scriptsize{[91,100]} & 100 & 13.4 & 1.8 & 24 \\}
\def\TabFailures{Unmet precondition & 25 & 16 & 3 & 2 \\
Wrong object kind & 1 & 1 & 0 & 0 \\
Wrong pour direction & 3 & 2 & 2 & 1 \\
Ends holding (advisory) & 5 & 3 & 2 & 1 \\
Pointless move (advisory) & 12 & 8 & 6 & 4 \\}
\def\TabSplit{GPT-4.1 mini & 93 & 83 \\
GPT-4.1 & 100 & 100 \\
Claude Sonnet 5.5 & 100 & 100 \\
Claude Opus 5.5 & 100 & 100 \\}
\def\TabGapData{EPIC-KITCHENS-100 & 84 \\
EgoDex & 55 \\
Egocentric-10K & 12 \\
Ego-Exo4D & 11 \\
Wipe whiteboard (villekuosmanen) & 10 \\}
\def\PlotRepairs{\addplot coordinates {(0,56.4) (1,84.6) (2,87.2)};
\addlegendentry{GPT-4.1 mini}
\addplot coordinates {(0,89.5) (1,97.4) (2,100.0)};
\addlegendentry{GPT-4.1}
\addplot coordinates {(0,87.2) (1,100.0) (2,100.0)};
\addlegendentry{Claude Sonnet 5.5}
\addplot coordinates {(0,94.9) (1,100.0) (2,100.0)};
\addlegendentry{Claude Opus 5.5}}
\def\TabRequests{Paint an image & N & \cellcolor{green!22}2/3 & \cellcolor{green!33}3/3 & \cellcolor{green!33}3/3 & \cellcolor{green!33}3/3 \\
Sandcastle & N & \cellcolor{green!4}0/3$^\dagger$ & \cellcolor{green!22}2/3 & \cellcolor{green!22}2/3 & \cellcolor{green!22}2/3 \\
Card trick & N & \cellcolor{green!4}0/3 & \cellcolor{green!33}3/3 & \cellcolor{green!33}3/3 & \cellcolor{green!33}3/3 \\
Water plants & N & \cellcolor{green!11}1/3$^\dagger$ & \cellcolor{green!4}0/3 & \cellcolor{green!33}3/3 & \cellcolor{green!33}3/3 \\
Write on whiteboard & N & \cellcolor{green!33}3/3 & \cellcolor{green!33}3/3 & \cellcolor{green!33}3/3 & \cellcolor{green!33}3/3 \\
Jam sandwich & N & \cellcolor{green!33}3/3 & \cellcolor{green!33}3/3 & \cellcolor{green!33}3/3 & \cellcolor{green!33}3/3 \\
Flowers in vase & N & \cellcolor{green!33}3/3 & \cellcolor{green!33}3/3 & \cellcolor{green!33}3/3 & \cellcolor{green!33}3/3 \\
Feed the cat & N & \cellcolor{green!33}3/3 & \cellcolor{green!33}3/3 & \cellcolor{green!22}2/3 & \cellcolor{green!33}3/3 \\
Pi\~na colada & D & \cellcolor{green!4}0/3 & \cellcolor{green!33}3/3 & \cellcolor{green!22}2/3 & \cellcolor{green!33}3/3 \\
Glass of juice & D & \cellcolor{green!33}3/3 & \cellcolor{green!33}3/3 & \cellcolor{green!33}3/3 & \cellcolor{green!33}3/3 \\
Smoothie & D & \cellcolor{green!11}1/3$^\dagger$ & \cellcolor{green!33}3/3 & \cellcolor{green!11}1/3 & \cellcolor{green!22}2/3 \\
Wipe table & D & \cellcolor{green!4}0/3 & \cellcolor{green!33}3/3 & \cellcolor{green!33}3/3 & \cellcolor{green!33}3/3 \\
Clear table & D & \cellcolor{green!33}3/3 & \cellcolor{green!33}2/2 & \cellcolor{green!33}3/3 & \cellcolor{green!33}3/3 \\}
\def\TabStability{GPT-4.1 mini & 0.86 & 0.74 & 23 & 0.84 \\
GPT-4.1 & 0.93 & 0.84 & 38 & 0.87 \\
Claude Sonnet 5.5 & 0.85 & 0.70 & 28 & 0.68 \\
Claude Opus 5.5 & 0.90 & 0.84 & 31 & 0.82 \\}
\def\TabVerbs{\code{drag} & 31 & 9 & \code{drag\_on\_surface}, \code{drag\_along\_surface} \\
\code{press} & 31 & 7 & \code{press\_into}, \code{press\_button} \\
\code{tilt} & 16 & 3 & \code{tilt\_out}, \code{tilt\_out\_food} \\
\code{dip} & 15 & 8 & \code{dip\_into}, \code{dip\_into\_vessel} \\
\code{scoop} & 15 & 7 & \code{scoop\_item}, \code{scoop\_material} \\
\code{peel} & 15 & 3 & \code{peel\_off}, \code{peel\_item} \\
\code{flip} & 15 & 3 & \code{flip\_item}, \code{flip\_over} \\
\code{twist} & 13 & 5 & \code{twist\_item}, \code{twist\_open} \\}
\def\TabBids{1 & \code{press\_to\_surface} & 10.41 & \code{press} & 20.21 \\
2 & \code{tilt\_out} & 7.39 & \code{drag} & 14.45 \\
3 & \code{run\_blender} & 6.54 & \code{tilt} & 10.74 \\
4 & \code{drag\_along\_surface} & 6.28 & \code{twist} & 10.40 \\
5 & \code{peel\_item} & 4.96 & \code{peel} & 9.94 \\
6 & \code{twist\_item} & 3.86 & \code{run\_blender} & 8.79 \\}
\def\TabCategories{Cleaning & 6 & 1 & 480 \\
Finishing & 5 & 0 & 398 \\
Household & 16 & 0 & 428 \\
Manipulation & 13 & 4 & 444 \\
Mixing & 5 & 1 & 309 \\
Perception & 5 & 2 & 950 \\
Pouring & 4 & 1 & 618 \\
Transport & 5 & 2 & 1{,}141 \\}
\def\PlotSpeedup{\addplot[only marks, mark=*, mark size=1.2pt, gray, fill opacity=0.3, draw opacity=0.4] coordinates {(10,5) (9,5) (10,5) (9,8) (8,8) (12,12) (12,12) (16,16) (5,5) (5,5) (4,4) (14,14) (11,7) (14,14) (10,10) (10,10) (10,10) (6,6) (6,6) (6,6) (26,26) (33,27) (26,26) (6,6) (6,6) (6,6) (24,24) (21,21) (5,5) (5,5) (5,5) (8,8) (8,8) (8,8) (6,6) (6,6) (6,6) (11,8) (11,7) (17,12) (4,4) (4,4) (4,4) (7,7) (9,9) (7,7) (5,5) (5,5) (5,5) (14,9) (14,10) (14,10) (10,10) (10,10) (10,10) (9,5) (9,5) (8,4) (24,24) (25,25) (24,24) (6,6) (6,6) (6,6) (17,17) (17,17) (21,15) (6,6) (6,6) (6,6) (20,20) (20,20) (5,5) (5,5) (7,7) (14,12) (11,11) (10,6) (7,5) (8,5) (6,4) (8,8) (6,6) (8,8) (5,5) (5,5) (5,5) (14,14) (17,12) (12,12) (10,10) (10,10) (10,10) (13,13) (11,7) (12,12) (26,26) (26,26) (26,26) (6,6) (6,6) (6,6) (21,21) (17,17) (17,17) (5,5) (5,5) (5,5) (8,8) (8,8) (8,8) (7,7) (7,7) (7,7) (11,7) (12,8) (13,9) (16,12) (8,5) (16,12) (8,8) (8,8) (8,8) (13,8) (13,9) (12,8) (18,12) (17,13) (19,19) (10,10) (10,10) (10,10) (14,10) (13,10) (14,10) (26,26) (26,26) (26,26) (6,6) (6,6) (6,6) (25,25) (25,25) (25,25) (5,5) (7,7) (5,5) (16,16) (20,20) (16,16)};
\addlegendentry{default hands}
\addplot[only marks, mark=triangle*, mark size=1.5pt, red!70!black, fill opacity=0.5, draw opacity=0.6] coordinates {(10,5) (9,5) (10,5) (9,5) (8,4) (12,8) (12,8) (16,10) (5,5) (5,5) (4,4) (14,8) (11,7) (14,9) (10,8) (10,8) (10,8) (6,6) (6,6) (6,6) (26,17) (33,18) (26,17) (6,6) (6,6) (6,6) (24,14) (21,12) (5,5) (5,5) (5,5) (8,4) (8,4) (8,4) (6,6) (6,6) (6,6) (11,8) (11,7) (17,12) (4,4) (4,4) (4,4) (7,7) (9,9) (7,7) (5,5) (5,5) (5,5) (14,9) (14,10) (14,10) (10,8) (10,8) (10,8) (9,5) (9,5) (8,4) (24,15) (25,16) (24,14) (6,6) (6,6) (6,6) (17,12) (17,12) (21,13) (6,6) (6,6) (6,6) (20,12) (20,12) (5,5) (5,5) (7,7) (14,8) (11,6) (10,6) (7,5) (8,5) (6,4) (8,8) (6,6) (8,8) (5,5) (5,5) (5,5) (14,8) (17,10) (12,8) (10,7) (10,8) (10,7) (13,9) (11,7) (12,6) (26,17) (26,17) (26,17) (6,6) (6,6) (6,6) (21,12) (17,12) (17,12) (5,5) (5,5) (5,5) (8,4) (8,4) (8,4) (7,7) (7,7) (7,7) (11,7) (12,8) (13,9) (16,12) (8,5) (16,12) (8,8) (8,8) (8,8) (13,8) (13,9) (12,8) (18,11) (17,10) (19,11) (10,8) (10,8) (10,8) (14,10) (13,9) (14,10) (26,17) (26,17) (26,17) (6,6) (6,6) (6,6) (25,15) (25,15) (25,15) (5,5) (7,7) (5,5) (16,8) (20,12) (16,8)};
\addlegendentry{balanced}}


\begin{document}

\twocolumn[{%
\begin{center}
{\LARGE\bfseries Act It Out, Then Bind: Verifier-in-the-Loop Compilation of\\[2pt]
Natural-Language Requests into Reusable Robot Skills\par}
\vspace{10pt}
{\large Frenchies}\\[2pt]
{\normalsize Z\"urich}\\[2pt]
\vspace{6pt}
\end{center}
\begin{center}
\begin{minipage}{0.86\textwidth}
\small
\textbf{Abstract.} A robot skill marketplace has to turn a customer's sentence
(``mix me a smoothie'') into a sequence of purchasable skills that a robot can
actually execute, and say precisely what is missing when it cannot. We present
\sys, which compiles requests in two stages: a language model first \emph{acts
out} the request as atomic hand movements without seeing the catalog, then
\emph{binds} those movements to catalog skills under a schema that only admits
real skill identifiers. A symbolic simulator over hand-scoped STRIPS-style
predicates, object kinds and a contents ledger checks every plan and returns its
complaints to the model for repair. Uncovered movements become atomic, reusable
\emph{gaps}, each matched to an open demonstration dataset, and unserved gaps are
priced as bids for data. On \NAll{} compilations of 13 requests by four planners
(GPT-4.1 mini, GPT-4.1, Claude Sonnet~5.5, Claude Opus~5.5), \BlockedFirstAll\% of
first bindings are rejected by the simulator and \BlockedFinalAll\% remain rejected
after two repair rounds; with feedback, mid-sized planners match the largest at a
fraction of its latency. Most first-binding failures are unmet preconditions, but
some pass every precondition check and are caught only by object kinds and the
contents ledger. Searching hand assignments shortens the resulting two-arm
schedules by \TimeSavedBal\%. \ShareGap\% of plan steps are gaps; nearly all map
to an open dataset, yet mostly to human egocentric video rather than robot
demonstrations, and their names must be grouped before demand can be priced.
\end{minipage}
\end{center}
\vspace{10pt}
}]

\section{Introduction}
Robot skills are becoming products: an operator buys \emph{grasp bottle} or
\emph{wipe surface} for a specific arm, and expects them to compose into a process.
Vision-language-action policies \cite{brohan2023rt2,kim2024openvla,octo2024,black2024pi0}
make such skills cheaper to train from a few dozen demonstrations, and open corpora
\cite{oxe2024,khazatsky2024droid,cadene2024lerobot} make the demonstrations easier
to find. What remains is the interface. A customer does not shop for
\code{move\_to\_station}; they ask for a smoothie, a clean table or a painted
picture, and expect the store to work out which skills that takes, which of them
exist, and what is missing.

Natural language makes a language model the obvious compiler---and an unreliable
one. Language models produce plans that read well but are physically incoherent:
pouring before grasping, closing a bottle grasp around a bucket, holding a blender
jar while pressing its start button, or pouring a glass into the blender instead of
the other way round \cite{valmeekam2023planbench,kambhampati2024llmmodulo}. For a
store this is worse than an error message, because an incoherent plan is also a
wrong bill, a wrong installation on a physical arm, and a wrong estimate of what
data to collect next.

We take the position of LLM-Modulo frameworks \cite{kambhampati2024llmmodulo}: the
model proposes, a sound checker decides. Our contributions are:
\begin{itemize}
\item a \textbf{two-stage compiler} that separates \emph{what a person's hands
  would do} from \emph{which product does it}, so that uncovered work surfaces as
  atomic, reusable gaps instead of being dropped or collapsed into one opaque
  ``draw image with brush'' step (\S\ref{sec:compile});
\item a \textbf{symbolic simulator} whose checks are pure functions of declared
  data---skill pre/postconditions, the object kinds a skill acts on, and a ledger of
  what each container holds---with no task-specific rules (\S\ref{sec:sim});
\item an \textbf{episode-level scheduler} that splits a checked plan across two
  arms without trading validity for speed (\S\ref{sec:sched});
\item a \textbf{demand-to-data bidder} that prices unserved gaps per demonstration
  episode (\S\ref{sec:bids});
\item an \textbf{evaluation} of four planners with per-round repair traces, a
  taxonomy of the failures caught, run-to-run stability, schedule speed-ups and the
  open-data coverage of the compiled plans (\S\ref{sec:exp}--\S\ref{sec:results}).
\end{itemize}

\section{Related Work}
\textbf{Language models as planners.} Language models can decompose instructions
into steps \cite{huang2022zeroshot}, ground them in skill affordances
\cite{ahn2022saycan}, close the loop with environment feedback
\cite{huang2022inner}, emit programs over robot APIs
\cite{liang2023code,singh2023progprompt}, or compose spatial value maps for
low-level motion \cite{huang2023voxposer}. Benchmarks show their plans are
frequently invalid when checked formally \cite{valmeekam2023planbench}.

\textbf{Planners with verifiers.} LLM+P translates problems to PDDL
\cite{mcdermott1998pddl} for a classical planner \cite{liu2023llmp}; LLM-Modulo
pairs the model with external critics \cite{kambhampati2024llmmodulo}; Text2Motion
checks geometric feasibility of skill sequences \cite{lin2023text2motion}; task and
motion planning integrates both levels \cite{garrett2021tamp}. We use the critic as
a \emph{repair signal} rather than a search engine, because the catalog is
open-ended and a plan must also name, precisely, what is missing. Our operators
follow STRIPS \cite{fikes1971strips}, extended with hand-scoped facts.

\textbf{Skills and policies.} Skills in our catalog are policies such as VLA models
fine-tuned from demonstrations \cite{brohan2023rt2,kim2024openvla,octo2024,black2024pi0},
including bimanual ones from low-cost teleoperation \cite{zhao2023aloha}.

\textbf{Demonstration data.} Robot demonstrations come from teleoperation corpora
\cite{oxe2024,khazatsky2024droid,cadene2024lerobot}; human demonstrations from
egocentric video \cite{grauman2022ego4d,damen2022epic,grauman2024egoexo,hoque2025egodex}.
We use both as candidate training sources for missing skills, and measure which of
them the compiled plans actually point to.

\section{Method}
\begin{figure*}[t]
\centering
\begin{tikzpicture}[font=\footnotesize, node distance=5mm,
  box/.style={draw, rounded corners=2pt, align=center, text width=27mm,
              minimum height=12mm, inner sep=2pt, fill=#1},
  box/.default=gray!6, arr/.style={-latex, thick}]
\node[box] (req) {Request\\``mix me a smoothie''};
\node[box=blue!6, right=of req] (act) {\textbf{1. Act it out}\\objects + kinds,\\phases $\to$ movements};
\node[box=blue!6, right=of act] (bind) {\textbf{2. Bind}\\catalog skill (enum)\\or \code{GAP} + dataset};
\node[box=orange!10, right=of bind] (sim) {\textbf{3. Simulate}\\predicates per hand,\\\code{acts\_on}, contents};
\node[box=green!8, right=12mm of sim] (out) {\textbf{4. Schedule}\\two arms,\\episode search};
\node[box, text width=92mm, minimum height=7mm, below=6mm of $(bind.south)!0.5!(sim.south)$] (cat)
  {Catalog: \NSkills{} skills, \NPredicates{} predicates, \NKinds{} kinds\,$\cdot$\,\NOpenDatasets{} open datasets};
\node[box=red!5, text width=40mm, minimum height=7mm, below=6mm of out] (bid) {Unserved gaps $\to$ data bids};
\draw[arr] (req) -- (act); \draw[arr] (act) -- (bind); \draw[arr] (bind) -- (sim);
\draw[arr] (sim) -- node[above, font=\scriptsize]{no issue} (out);
\draw[arr, red!70!black] (sim.north) to[bend right=30] node[above, font=\scriptsize]{issues + catalog hints ($\le$2 rounds)} (bind.north);
\draw[arr, gray] (cat.north -| bind) -- (bind.south);
\draw[arr, gray] (cat.north -| sim) -- (sim.south);
\draw[arr, gray] (out) -- (bid);
\end{tikzpicture}
\caption{\sys{} compiles a request in two model calls, checks the result with a
declarative simulator whose complaints, with catalog-derived hints, go back to the
binder, and schedules the valid plan on two arms. Only the binder sees the catalog.
Plans that cannot run feed the data bidder.}
\label{fig:pipeline}
\end{figure*}

\subsection{Skills as operators}
The catalog holds \NSkills{} skills (\NTrained{} trained, \NBeta{} beta,
\NCatalogOnly{} listed without a policy). Each declares preconditions,
postconditions and deletes over \NPredicates{} predicates. Ten predicates are
\emph{hand-scoped} (\code{holding}, \code{hand\_empty}, \code{object\_located},
\ldots): \code{holding@left} and \code{holding@right} are distinct facts, which is
what lets two sub-chains run on two arms. \NActsOn{} skills declare which of
\NKinds{} object kinds they \emph{act on} (a bottle grasp acts on \code{bottle}, an
item grasp on \code{item}); pouring skills declare a flow (held $\rightarrow$
target). Kinds are coarse on purpose---\code{bottle}, \code{glass},
\code{shaker}, \code{scoop}, \code{cloth}, \code{item}, \code{fixed},
\code{substance}---because they describe what a gripper closes around, not what an
object is called; the full list is in Appendix~\ref{app:requests}.

\subsection{Two-stage compilation}\label{sec:compile}
\textbf{Act it out.} Without the catalog in view, the model writes the request as
(i)~the objects involved, each classified into an object kind, (ii)~phases named by
their outcome (``the glass holds juice''), each broken into movements that start
with a movement verb (pick up, carry, tilt, press, drag), and (iii)~a goal. Outcome
verbs (make, operate, prepare) may name a phase but never a movement. A short
start-state description fixes everyday assumptions (bottles open, appliances on
their base and operated with a free hand). Keeping the catalog out of this stage
matters: a model that plans \emph{from} the catalog bends the task to fit it, and
whatever the catalog cannot do is dropped or squashed into one gap the size of the
request.

\textbf{Bind.} With the catalog and the open-dataset list in view, a second call
maps movements to skills or to \code{GAP} entries. Structured outputs restrict skill
identifiers to an enum of real skills, targets to the object names from stage one,
and dataset identifiers to open datasets, so a hallucinated skill is impossible by
construction. A gap is named as a reusable motion primitive (\code{twist\_open},
\code{press\_to\_surface}) with conditions in the catalog's vocabulary; its
position in the plan is exact because gaps are emitted inline with the steps.
Fig.~\ref{fig:example} shows a compiled plan.

\begin{figure}[t]
\centering\footnotesize
\setlength{\tabcolsep}{3pt}
\begin{tabular}{@{}rl@{}}
\toprule
\multicolumn{2}{@{}l}{\emph{``You have a paint bucket and brush. Draw an image on the table.''}}\\
\multicolumn{2}{@{}l}{Objects: paint bucket \code{item}, paint \code{substance},}\\
\multicolumn{2}{@{}l}{\phantom{Objects: }paintbrush \code{item}, table \code{fixed}}\\
\midrule
1 & \code{locate\_object}(paintbrush)\\
2 & \code{grasp\_object}(paintbrush)\\
3 & \code{move\_to\_station}(paint bucket)\\
  & \textcolor{red!70!black}{\code{GAP dip\_into}} $\to$ EPIC-KITCHENS-100\\
4 & \code{move\_to\_station}(paint bucket)\\
5 & \code{move\_to\_station}(table)\\
  & \textcolor{red!70!black}{\code{GAP drag\_along\_surface}} $\to$ EPIC-KITCHENS-100\\
6 & \code{move\_to\_station}(paint bucket)\\
7 & \code{release\_object}\\
\midrule
\multicolumn{2}{@{}l}{Goal: an image is painted on the table and both hands are empty.}\\
\bottomrule
\end{tabular}
\caption{A compiled plan, verbatim from the Claude Opus~5.5 run (valid on the first
binding). The grasp is an \emph{item} grasp because stage one classified the brush
as an item; the painting itself is two reusable gaps rather than one ``draw image''
step. Step~4 repeats step~3, an inefficiency the simulator does not flag because the
brush is used in between.}
\label{fig:example}
\end{figure}

\subsection{Simulation and repair}\label{sec:sim}
The simulator executes the plan step by step, gaps included with the effects they
declare, and assigns hands by grouping steps into \emph{episodes} (from the first
step that leaves something on a hand to the step that frees it); overlapping
episodes go to opposite hands. It raises \emph{blocking} issues---unmet
precondition, wrong object kind (\code{acts\_on} $\not\ni$ kind of target), wrong
pour direction (from a ledger of which object each hand holds and what each
container contains), embodiment mismatch---and \emph{advisory} ones: a plan ending
with a hand still holding something, or a \emph{pointless move} (an object picked
up and set back where it was without being used). Issues are returned to the binder
with hints derived from the catalog (``\code{positioned} is established by
\code{move\_to\_station}''; ``what holds something now: blender''). The loop runs
at most two rounds (plus one for a pour in the wrong direction), and only the
binding is retried: the acted-out script is not in question. Every number in this
paper counts a plan as \emph{valid} only if the simulator raises no blocking issue.
Appendix~\ref{app:checks} lists every check.

\subsection{Two-arm scheduling}\label{sec:sched}
A valid plan is turned into a dependency graph with four kinds of edge, all
physical: \emph{requires} (a predicate an earlier step established),
\emph{consumes} (clearing a predicate an earlier step established), \emph{hand}
(the same hand is still busy) and \emph{resource} (two steps write the same world
fact, or two arms want the same place---only one bottle fits over the glass). Steps
are layered by their longest dependency chain; the number of layers is the
schedule's depth. Hands are chosen per episode, not per step, because a hand has
to finish what it started: we search all assignments exactly up to 13 episodes
and hill-climb beyond, pin the first episode to the left hand (mirror schedules
are equivalent), and reject any assignment that adds a blocking issue.

\subsection{From gaps to data bids}\label{sec:bids}
Each request $r$ whose plan cannot run is logged against its blockers $B_r$
(gaps and untrained skills). For a blocker $s$, the $k$-th identical prompt counts
$1/k$, and requests fade with age $a_r$ at half-life $h$:
\begin{equation}
w_{r,s} = \frac{2^{-a_r/h}}{k_{r,s}}, \qquad
V_s = \sum_{r \ni s} w_{r,s}\,\frac{P_r}{|B_r|},
\end{equation}
where $P_r$ is the price of the whole plan (catalog prices, each gap at the
catalog's median price), shared equally among its blockers because all of them must
be filled. With $H$ future sales per logged request, target return $\rho$, a cap
$c$ and $E_s$ episodes needed ($E_0$, reduced by a factor $1-\delta$ when an open
dataset matches), the bid per episode and the budget pacing factor are
\begin{equation}
b_s = \min\!\Big(c,\ \frac{H\,V_s}{\rho\,E_s}\Big), \qquad
\alpha = \min\!\Big(1,\ \frac{\text{budget}}{\sum_s b_s E_s}\Big),
\end{equation}
and the posted bid is $\alpha\, b_s$. Pacing scales all bids equally, so more unmet
demand always means a higher bid. Defaults: budget 5{,}000\,EUR per month,
$\rho=3$, $H=10$, $E_0=50$, $\delta=0.4$, $h=30$ days, $c=60$\,EUR.

\section{Experimental Setup}\label{sec:exp}
\textbf{Requests.} Five demo requests (pi\~na colada, juice, smoothie, wiping,
clearing up) and eight unfamiliar ones (painting with a brush, a sandcastle, a card
trick, watering plants, writing on a whiteboard, a jam sandwich, flowers in a vase,
feeding a cat); the exact wording is in Appendix~\ref{app:requests}.
\textbf{Planners.} GPT-4.1 mini, GPT-4.1, Claude Sonnet~5.5 and Claude Opus~5.5
(effort \emph{low}, standard speed), the same model for both stages and for repair.
\textbf{Protocol.} Each request is compiled \Runs{} times per planner with the
compile cache off; the compiler records the simulator's issues at the start of
every repair round, so a single run yields validity after 0, 1 and 2 rounds. We
report proportions with 95\% Wilson intervals \cite{wilson1927}; \ApiFailures{}
compilation that failed on a provider rate limit is excluded.

\textbf{Data.} \NDatasets{} datasets are registered (\NOpenDatasets{} openly
accessible: \NTeleop{} robot teleoperation, \NEgo{} egocentric human video,
\NSim{} simulation). \NLinks{} curated links connect each catalog skill to
teleoperation datasets whose task text matches it, \TotalEpisodes{} episodes in
total (median \MedianEpisodes{} per skill, minimum \MinEpisodes;
Table~\ref{tab:categories}). Every link was checked for licence, access, episode
count and task text.

\begin{table}[t]
\centering\small
\caption{Catalog by category: skills, trained skills, and the median number of
curated open teleoperation episodes per skill.}
\label{tab:categories}
\begin{tabular}{lrrr}
\toprule
Category & Skills & Trained & Median episodes \\
\midrule
\TabCategories
\bottomrule
\end{tabular}
\end{table}

\section{Results}\label{sec:results}

\begin{table*}[t]
\centering\small
\setlength{\tabcolsep}{6pt}
\caption{Planner comparison over 13 requests $\times$ \Runs{} runs. Valid after
$k$ repair rounds (\%), final validity with 95\% Wilson interval, plans ending with
both hands empty (\%), mean steps and gaps per plan, median seconds per compile.}
\label{tab:models}
\begin{tabular}{lccccccc}
\toprule
Planner & $k{=}0$ & $k{=}1$ & Final & Empty & Steps & Gaps & s \\
\midrule
\TabModels
\bottomrule
\end{tabular}
\end{table*}

\begin{figure}[t]
\centering
\begin{tikzpicture}
\begin{axis}[width=\columnwidth, height=4.4cm,
  xlabel={repair rounds}, ylabel={valid plans (\%)},
  xtick={0,1,2}, xticklabels={0,1,2 (final)},
  ymin=40, ymax=102, grid=major, grid style={gray!25},
  legend style={font=\scriptsize, at={(0.98,0.03)}, anchor=south east},
  every axis plot/.append style={thick, mark size=1.8pt},
  cycle list={{blue!70!black,mark=*},{teal,mark=square*},{orange!90!black,mark=triangle*},{red!70!black,mark=diamond*}}]
\PlotRepairs
\end{axis}
\end{tikzpicture}
\caption{Share of plans the simulator accepts after each round of feedback.}
\label{fig:repairs}
\end{figure}

\subsection{Simulator feedback does most of the work}
Without feedback the planners differ widely: \ValidZeroGPTmini\% of GPT-4.1 mini's
first bindings pass the simulator, against \ValidZeroOpus\% for Claude Opus~5.5
(Table~\ref{tab:models}, Fig.~\ref{fig:repairs}). One round of repair closes most of
that difference: GPT-4.1, Sonnet~5.5 and Opus~5.5 reach \ValidOneGPT\%,
\ValidOneSonnet\% and \ValidOneOpus\%, and GPT-4.1 mini rises to
\ValidOneGPTmini\%. The second round adds little. With the checker in the loop,
GPT-4.1 and Sonnet~5.5 match Opus~5.5's final validity at a fraction of its latency
(median \SecsGPT{}\,s and \SecsSonnet{}\,s against \SecsOpus{}\,s per compilation).
Beyond the smallest planner, the checker rather than model scale makes the plans
installable; scale mostly buys fewer repair calls.

\begin{table}[t]
\centering\small
\caption{Issues raised at the first binding and on the final plan, pooled over all
\NAll{} compilations (plans affected, \%).}
\label{tab:failures}
\begin{tabular}{lrrrr}
\toprule
& \multicolumn{2}{c}{First binding} & \multicolumn{2}{c}{Final} \\
Issue & $n$ & \% & $n$ & \% \\
\midrule
\TabFailures
\bottomrule
\end{tabular}
\end{table}

\subsection{What the simulator catches}
Of the \IssuesFirst{} blocking issues raised on first bindings, \UnmetFirst{} are
unmet preconditions (Table~\ref{tab:failures}), spread across requests---for
instance a locate on a hand that is still holding something, or a release without a
preceding carry. The remaining \FlowFirst{} wrong pour directions and \ObjectFirst{}
wrong object kind satisfy every predicate and would pass a precondition-only
checker; they are caught only by the contents ledger and the \code{acts\_on}
declarations. After repair, \BlockedFinalN{} of \NAll{} plans stay blocked, all from
GPT-4.1 mini (a sandcastle, watering plants by pouring from an empty glass, and a
smoothie). Advisory \emph{pointless moves} halve with feedback
(\PointlessFirst{}$\rightarrow$\PointlessFinal) but persist, as the binder
sometimes keeps a carry the first stage asked for.

\begin{table}[t]
\centering\small
\setlength{\tabcolsep}{4pt}
\caption{First bindings the simulator accepts, per request and planner (out of
\Runs; shading by share). $^\dagger$: not every run valid after repair. N: unfamiliar,
D: demo request.}
\label{tab:requests}
\begin{tabular}{llcccc}
\toprule
Request & & 4.1 mini & 4.1 & Sonnet & Opus \\
\midrule
\TabRequests
\bottomrule
\end{tabular}
\end{table}

\subsection{Which requests are hard}
Difficulty is concentrated (Table~\ref{tab:requests}). Writing on a whiteboard,
the jam sandwich, flowers in a vase and a glass of juice pass on the first binding
for every planner and run. The sandcastle---two tools, a substance and a tray---is
the only request no planner gets right every time on the first binding, and the
smoothie trips three of four planners at least once. GPT-4.1 mini fails the first binding
of the card trick, the pi\~na colada and the table wiping in every run, yet repair
fixes all of them; its remaining failures are the requests that need two tools or
a pour.

\begin{table}[t]
\centering\small
\setlength{\tabcolsep}{4pt}
\caption{Stability across the \Runs{} runs of a request: mean pairwise similarity
(1 $-$ normalised edit distance) of skill sequences, with targets, share of
identical plans, and Jaccard overlap of gap verbs.}
\label{tab:stability}
\begin{tabular}{lcccc}
\toprule
Planner & Skills & +Targets & Identical (\%) & Gaps \\
\midrule
\TabStability
\bottomrule
\end{tabular}
\end{table}

\subsection{Plans are similar, not identical}
Runs of the same request produce similar plans---mean skill-sequence similarity
between \StabSonnet{} and \StabGPT{} (Table~\ref{tab:stability})---but rarely the
same plan: only \SameGPTmini--\SameGPT\% of run pairs are identical. Differences are
mostly in targets (which object is carried where) and in how a missing motion is
split into gaps. For a store, this means a price quoted for a request is a property
of one compilation, not of the request; caching compiled plans, as \sys{}
does, makes quotes reproducible.

\subsection{Two arms help once hands are searched}
With the simulator's default hand assignment---each new episode on the first free
hand---only \ParShare\% of the \NScheduled{} valid plans run anything in parallel
(mean speed-up \SpeedMeanDefault$\times$, Fig.~\ref{fig:speedup}): the compiler is
told to use one hand at a time unless two are needed, which keeps plans valid but
serial. Searching hand assignments per episode raises this to \ParShareBal\% of
plans, a mean speed-up of \SpeedMeanBal$\times$ (up to \SpeedMaxBal$\times$), and
\TimeSavedBal\% fewer time slots overall, without adding a single blocking issue.
Parallelism is available in most household requests; it has to be found after
validation, not asked of the model.

\begin{figure}[t]
\centering
\begin{tikzpicture}
\begin{axis}[width=\columnwidth, height=4.6cm,
  xlabel={plan length (steps)}, ylabel={schedule depth},
  xmin=0, xmax=36, ymin=0, ymax=30, grid=major, grid style={gray!25},
  legend style={font=\scriptsize, at={(0.03,0.97)}, anchor=north west}]
\PlotSpeedup
\addplot[dashed, black, domain=0:30] {x};
\addlegendentry{serial ($y{=}x$)}
\end{axis}
\end{tikzpicture}
\caption{Schedule depth against plan length for all valid plans, with the default
hand assignment and after the episode search. Points below the diagonal run steps
on both arms at once.}
\label{fig:speedup}
\end{figure}

\subsection{Open-data coverage}
Across all final plans, \ShareTrained\% of steps bind to trained skills,
\ShareUntrained\% to listed but untrained skills, and \ShareGap\% are gaps. Of
\GapTotal{} gaps, \GapWithData\% were matched to an open dataset---but
\GapEgoShare\% of those matches point to human egocentric video
(Table~\ref{tab:gapdata}) rather than robot teleoperation. The motions the catalog
lacks (dipping a brush, spreading, scooping, twisting a lid open) are exactly those
without robot demonstrations; human video can pretrain them, but not replace robot
episodes. The catalog's own skills, by contrast, are well covered by teleoperation:
every skill has at least \MinEpisodes{} curated episodes (Table~\ref{tab:categories}).

\begin{table}[t]
\centering\small
\caption{Open datasets most often proposed for gaps (gap count).}
\label{tab:gapdata}
\begin{tabular}{lr}
\toprule
Dataset & Gaps \\
\midrule
\TabGapData
\bottomrule
\end{tabular}
\end{table}

\subsection{Demand needs canonical names}
Gap names are not canonical: the \GapTotal{} gaps carry \DistinctGaps{} distinct
names, and dipping a tool into a container alone appears under eight
(Table~\ref{tab:verbs}). Grouping names by their leading verb collapses them to
\DistinctVerbs{} motions. This matters for the bidder: treating each of the
\NAll{} compilations as a customer request, the market splits demand over
\MarketSkillsRaw{} distinct blockers with raw names but \MarketSkillsCanon{} after
grouping, and the top bid per episode rises from \TopBidRaw{} to \TopBidCanon{}\,EUR
(Table~\ref{tab:bids}): fragmented names hide the most wanted motions. Verb grouping
is a blunt instrument---it merges pressing a button with pressing a brush onto a
surface---so a deployed store needs a proper motion taxonomy.

\begin{table}[t]
\centering\small
\setlength{\tabcolsep}{3pt}
\caption{Most frequent gap verbs: gaps, distinct names, examples.}
\label{tab:verbs}
\begin{tabular}{@{}lrr>{\raggedright\arraybackslash}p{0.5\columnwidth}@{}}
\toprule
Verb & Gaps & Names & Examples \\
\midrule
\TabVerbs
\bottomrule
\end{tabular}
\end{table}

\begin{table}[t]
\centering\small
\setlength{\tabcolsep}{3pt}
\caption{Simulated market: top bids per episode (EUR) with raw gap names and with
names grouped by verb. Every compilation counts as one request; the default
parameters of \S\ref{sec:bids} apply.}
\label{tab:bids}
\footnotesize
\begin{tabular}{rlrlr}
\toprule
& Raw name & Bid & Grouped & Bid \\
\midrule
\TabBids
\bottomrule
\end{tabular}
\end{table}

\section{Discussion}
\textbf{A checker substitutes for scale.} The largest planner is the most reliable
on its own, but with two rounds of simulator feedback, planners that are roughly three
times faster reach the same final validity. For an interactive store this is
the right trade: the checker is cheap and exact, the model is slow and approximate,
and the store pays the model's latency on every request.

\textbf{The data gap is the dexterity gap.} The catalog's transport and grasping
skills have hundreds to tens of thousands of open robot episodes; the motions that
compiled plans actually miss---dipping, spreading, peeling, twisting---have almost
none and are matched to human video instead. Open robot data is broad in tasks but
shallow in dexterous tool use, which is where demand concentrates.

\textbf{Safety is a property of the pipeline.} Because the compiler can be wrong,
the gates are placed after it: installation and dispatch re-run the simulator on the
current catalog rather than trusting a stored verdict, a queued plan cannot be
edited, a runner claims a job atomically so two arms never receive the same plan, a
plan whose runner died is failed rather than silently re-run, and an operator stop
never sends the arm home. None of these depend on the model behaving well.

\textbf{From valid to certified.} A store should eventually publish, per skill and
per robot, a measured success rate with a confidence interval; with 20 trials,
19/20 is 95\% with a Wilson interval of roughly 76--99\%, which is why our runner
logs every attempt, failures included. Validity is the necessary first gate;
measured success is the product.

\section{Limitations}\label{sec:limits}
\textbf{Valid is not successful.} The simulator checks what the catalog declares;
a plan it accepts can still fail on hardware, and skills' success rates in our
system so far come from a mock executor, so we report none. \textbf{The requests
are not a held-out test set.} The prompts and start-state assumptions were
revised while inspecting failures on these same requests (e.g.\ the appliance
assumption after a smoothie plan held the blender jar while running it), so the
numbers measure the final system on its development set. \textbf{Small sample.}
Thirteen requests and \Runs{} runs per planner give wide intervals; differences of a
few points between planners are not significant. \textbf{Matching is declared, not
demonstrated.} Dataset links rest on task-text matches and, for gaps, on the model's
choice; we have not trained skills from them. \textbf{The market is simulated.}
Bids in Table~\ref{tab:bids} come from our own compilations, where the same request
recurs twelve times and is damped accordingly; real demand will differ.
\textbf{Scheduling ignores gaps' timing.} Speed-ups are computed over catalog steps
only, since gaps have no duration until they are trained.

\section{Conclusion}
Separating \emph{acting out} a request from \emph{binding} it to products, and
letting a declarative simulator veto and explain, turns language models into
usable compilers for a skill store: plans become checkable, schedules exploit both
arms once hands are searched, gaps become atomic products with candidate data, and
failures become named, countable categories. The next steps are closing the loop on
hardware---replacing ``valid'' with measured success---canonicalising gap names into
a motion taxonomy, and evaluating on requests the system has never seen.

{\small
\bibliographystyle{abbrvnat}
\bibliography{refs}
}

\appendix
\section{Requests and object kinds}\label{app:requests}
\textbf{Unfamiliar requests (N).} ``You have a paint bucket and brush. draw an image
on the table''; ``build a small sandcastle on the tray with the bucket and shovel'';
``do a simple card trick with the deck on the table''; ``water the two plants on the
shelf''; ``write my name on the whiteboard with the marker''; ``make me a jam
sandwich''; ``put the flowers in the vase and add water''; ``feed the cat: open the
tin and put the food in her bowl''. \textbf{Demo requests (D).} ``make me a pina
colada''; ``pour me a glass of juice''; ``mix me a smoothie''; ``wipe my table clean
after dinner''; ``clear the table and put the bottles away''.

\textbf{Object kinds.} \code{bottle}: a rigid bottle, carton or jug of liquid, held
upright; \code{glass}: a drinking glass, cup or mug; \code{shaker}: a cocktail
shaker; \code{scoop}: a measuring scoop with a handle; \code{cloth}: something soft
and deformable; \code{item}: any other rigid, hand-sized loose thing; \code{fixed}:
something that is not picked up (a table, a shelf, an appliance); \code{substance}:
a material inside something else (paint, water, sand).

\section{Simulator checks}\label{app:checks}
\noindent\begin{minipage}{\columnwidth}
\centering\footnotesize
\setlength{\tabcolsep}{3pt}
\begin{tabular}{@{}>{\raggedright\arraybackslash}p{0.31\columnwidth}>{\raggedright\arraybackslash}p{0.64\columnwidth}@{}}
\toprule
Check & Raised when \\
\midrule
Unmet precondition & a step needs a fact, on its hand, that no earlier step or gap established \\
Wrong object kind & the skill's \code{acts\_on} does not include the kind of its target \\
Wrong pour direction & the held object is the target, or is an empty glass or shaker \\
Embodiment mismatch & the skill needs a body (torso, base) the robot lacks \\
Unknown skill & a step names a skill outside the catalog (prevented by the enum) \\
Empty plan & the binding contains no executable step \\
\midrule
Ends holding \emph{(adv.)} & a hand still holds or has located something at the end \\
Pointless move \emph{(adv.)} & an object is picked up and set back where it was, unused \\
\bottomrule
\end{tabular}
\captionof{table}{Every check the simulator runs; \emph{adv.}: advisory, returned to
the binder but not blocking installation.}
\label{tab:checks}
\end{minipage}

\vspace{4pt}
\noindent\textbf{Repair hints.} Each issue returns with a catalog-derived hint: the
skills that establish a missing fact, the containers holding something, or the
skill a duplicate gap restates.

\end{document}
